\documentclass[11pt,a4paper]{ctexart}
\usepackage[margin=1.55cm]{geometry}
\usepackage{amsmath,amssymb}
\usepackage{booktabs,array,longtable}
\usepackage{xcolor}
\usepackage{tcolorbox}
\usepackage{hyperref}

\hypersetup{colorlinks=true, linkcolor=black, urlcolor=blue}
\setlength{\parindent}{0pt}
\setlength{\parskip}{5pt}
\renewcommand{\arraystretch}{1.35}
\newcommand{\dd}{\mathrm{d}}

\title{\vspace{-1.3em}粘性流动作业-2：题目与答案}
\author{}
\date{}

\begin{document}
\maketitle
\vspace{-2.2em}

\begin{tcolorbox}[colback=blue!3,colframe=blue!50!black,title=说明]
本文件按作业顺序整理，共 8 题。每题先列题目，再给答案。计算中取水密度未注明时为
\(\rho=1000\ \mathrm{kg/m^3}\)，重力加速度 \(g=9.81\ \mathrm{m/s^2}\)。
\end{tcolorbox}

\section*{1. 圆管湍流近壁分层}

\subsection*{题目}

直径为 \(50\ \mathrm{mm}\) 的管道内，流动的管截面平均的时均速度
\(\bar u_m\) 为 \(1.7\ \mathrm{m/s}\)，已知运动粘性系数
\[
\nu=1.006\ \mathrm{mm^2/s},
\]
计算壁面的线性底层、缓冲过渡层和完全湍流的厚度。提示：
\[
\frac{\bar u_m}{u_\tau}=2.5\ln\frac{Ru_\tau}{\nu}+1.75,
\]
\(R\) 是管半径。

\subsection*{解答}

\[
d=0.05\ \mathrm{m},\qquad R=\frac d2=0.025\ \mathrm{m}.
\]

\[
\nu=1.006\times10^{-6}\ \mathrm{m^2/s}.
\]

由平均速度公式反求摩擦速度 \(u_\tau\)：
\[
\frac{1.7}{u_\tau}
=2.5\ln\frac{0.025u_\tau}{1.006\times10^{-6}}+1.75.
\]

解得
\[
u_\tau\approx0.0818\ \mathrm{m/s}.
\]

近壁分层用
\[
y^+=\frac{yu_\tau}{\nu}.
\]

线性底层边界取 \(y^+=5\)：
\[
y_1=\frac{5\nu}{u_\tau}
=\frac{5\times1.006\times10^{-6}}{0.0818}
=0.0615\ \mathrm{mm}.
\]

缓冲层到完全湍流区的边界取 \(y^+=30\)：
\[
y_2=\frac{30\nu}{u_\tau}
=0.369\ \mathrm{mm}.
\]

所以从管壁向管心量距离 \(y\)：
\[
\boxed{0<y<0.0615\ \mathrm{mm}\quad\text{线性底层}}
\]

\[
\boxed{0.0615<y<0.369\ \mathrm{mm}\quad\text{缓冲过渡层}}
\]

\[
\boxed{0.369<y<25\ \mathrm{mm}\quad\text{完全湍流区}}
\]

若按“厚度”写：
\[
\boxed{\delta_{\text{线性底层}}=0.0615\ \mathrm{mm}}
\]
\[
\boxed{\delta_{\text{缓冲层}}=0.3075\ \mathrm{mm}}
\]
\[
\boxed{\delta_{\text{完全湍流区}}=24.63\ \mathrm{mm}}
\]

\section*{2. 泵吸水管最大安装高度}

\subsection*{题目}

如图离心泵的吸水管采用镀锌钢管，粗糙度 \(e=0.15\ \mathrm{mm}\)，长 \(6\ \mathrm{m}\)，直径 \(100\ \mathrm{mm}\)，有一个弯头
（局部阻力系数 \(\zeta=0.29\)）和一个进口阀（\(\zeta=2.5\)）。工作流量
\[
Q=50\ \mathrm{m^3/h},
\]
对应水温的汽化压力为 \(19.6\ \mathrm{kPa}\)，大气压力为 \(101.3\ \mathrm{kPa}\)。若要求水泵入口压力至少比汽化压力大 \(20\ \mathrm{kPa}\)，求水泵中心位置高于水面的最大安装高度 \(h\)。已知
\[
\nu=1.0\times10^{-6}\ \mathrm{m^2/s}.
\]

\subsection*{解答}

选水池自由液面为 A，泵入口中心为 B。取水面为零高度：
\[
z_A=0,\qquad z_B=h.
\]

水面速度近似为零：
\[
V_A\approx0.
\]

水面压强为大气压：
\[
p_A=p_a.
\]

泵入口最低允许压强：
\[
p_B=p_v+20\ \mathrm{kPa}=19.6+20=39.6\ \mathrm{kPa}.
\]

对 A 到 B 写机械能方程：
\[
p_a=p_B+\rho gh+\frac12\rho V^2+\Delta p_f+\sum\Delta p_\zeta.
\]

化为水头形式：
\[
h=\frac{p_a-p_B}{\rho g}
-\frac{V^2}{2g}
-\left(\lambda\frac{L}{d}+\sum\zeta\right)\frac{V^2}{2g}.
\]

流量：
\[
Q=\frac{50}{3600}=0.01389\ \mathrm{m^3/s}.
\]

管内平均速度：
\[
V=\frac{Q}{\pi d^2/4}
=\frac{0.01389}{\pi(0.1)^2/4}
=1.77\ \mathrm{m/s}.
\]

雷诺数：
\[
Re=\frac{Vd}{\nu}
=\frac{1.77\times0.1}{1.0\times10^{-6}}
=1.77\times10^5.
\]

相对粗糙度：
\[
\frac{e}{d}=\frac{0.15\ \mathrm{mm}}{100\ \mathrm{mm}}=0.0015.
\]

由 \(Re\) 和 \(e/d\) 查莫迪图，取
\[
\lambda\approx0.023.
\]

局部阻力系数和：
\[
\sum\zeta=0.29+2.5=2.79.
\]

\[
\frac{V^2}{2g}=\frac{1.77^2}{2\times9.81}=0.159\ \mathrm{m}.
\]

\[
\lambda\frac{L}{d}=0.023\times\frac{6}{0.1}=1.38.
\]

\[
\frac{p_a-p_B}{\rho g}
=\frac{(101.3-39.6)\times10^3}{1000\times9.81}
=6.29\ \mathrm{m}.
\]

因此：
\[
h=6.29-0.159-(1.38+2.79)\times0.159.
\]

\[
\boxed{h_{\max}\approx5.47\ \mathrm{m}}
\]

\section*{3. 水轮机输出功率}

\subsection*{题目}

如图，管路系统中有一水轮机，水温 \(10^\circ\mathrm{C}\)，对应运动粘度
\[
\nu=1.3\times10^{-6}\ \mathrm{m^2/s}.
\]
水管为铸铁管，粗糙度 \(e=0.25\ \mathrm{mm}\)。忽略局部损失，斜管中水流平均流速 \(2\ \mathrm{m/s}\)，试求水轮机输出功率。

图示条件：斜管内径 \(20\ \mathrm{cm}\)、长度 \(30\ \mathrm{m}\)；水平管内径 \(30\ \mathrm{cm}\)、长度 \(15\ \mathrm{m}\)；压力表读数 \(70\ \mathrm{kPa}\)；上游水面比压力表截面高 \(30\ \mathrm{m}\)。

\subsection*{解答}

斜管：
\[
d_1=0.20\ \mathrm{m},\quad L_1=30\ \mathrm{m},\quad v_1=2\ \mathrm{m/s}.
\]

水平管：
\[
d_2=0.30\ \mathrm{m},\quad L_2=15\ \mathrm{m}.
\]

流量：
\[
Q=\frac{\pi d_1^2}{4}v_1
=\frac{\pi(0.20)^2}{4}\times2
=0.0628\ \mathrm{m^3/s}.
\]

水平管速度：
\[
v_2=\frac{Q}{\pi d_2^2/4}
=\frac{0.0628}{\pi(0.30)^2/4}
=0.889\ \mathrm{m/s}.
\]

雷诺数：
\[
Re_1=\frac{v_1d_1}{\nu}
=\frac{2\times0.20}{1.3\times10^{-6}}
=3.08\times10^5.
\]

\[
Re_2=\frac{v_2d_2}{\nu}
=\frac{0.889\times0.30}{1.3\times10^{-6}}
=2.05\times10^5.
\]

相对粗糙度：
\[
\frac{e}{d_1}=0.00125,\qquad
\frac{e}{d_2}=0.000833.
\]

查莫迪图：
\[
\lambda_1\approx0.0216,\qquad
\lambda_2\approx0.0203.
\]

沿程损失：
\[
h_{f1}
=0.0216\times\frac{30}{0.20}\times\frac{2^2}{2\times9.81}
=0.659\ \mathrm{m}.
\]

\[
h_{f2}
=0.0203\times\frac{15}{0.30}\times\frac{0.889^2}{2\times9.81}
=0.041\ \mathrm{m}.
\]

\[
h_f=0.659+0.041=0.700\ \mathrm{m}.
\]

对上游水面 1 到压力表截面 2 写机械能方程：
\[
z_1=z_2+\frac{p_2}{\rho g}+\frac{v_2^2}{2g}+h_f+H_T.
\]

所以水轮机取走的水头：
\[
H_T=z_1-z_2-\frac{p_2}{\rho g}-\frac{v_2^2}{2g}-h_f.
\]

代入：
\[
H_T=30-\frac{70000}{1000\times9.81}
-\frac{0.889^2}{2\times9.81}
-0.700.
\]

\[
H_T=22.12\ \mathrm{m}.
\]

水轮机输出功率：
\[
P=\rho gQH_T.
\]

\[
P=1000\times9.81\times0.0628\times22.12.
\]

\[
\boxed{P\approx13.6\ \mathrm{kW}}
\]

\section*{4. 书 10.2：平板层流边界层}

\subsection*{题目}

\[
\mu=0.731\ \mathrm{Pa\cdot s},\qquad \rho=925\ \mathrm{kg/m^3}
\]
的油，以 \(0.6\ \mathrm{m/s}\) 的速度平行地流过一块长为 \(0.5\ \mathrm{m}\)、宽为 \(0.15\ \mathrm{m}\) 的光滑平板。求边界层最大厚度及平板所受的阻力。

\subsection*{解答}

运动粘度：
\[
\nu=\frac{\mu}{\rho}=\frac{0.731}{925}
=7.90\times10^{-4}\ \mathrm{m^2/s}.
\]

雷诺数：
\[
Re_L=\frac{UL}{\nu}
=\frac{0.6\times0.5}{7.90\times10^{-4}}
=3.80\times10^2.
\]

为层流。

边界层最大厚度：
\[
\delta_{\max}=\frac{5L}{\sqrt{Re_L}}
=\frac{5\times0.5}{\sqrt{3.80\times10^2}}
=0.128\ \mathrm{m}.
\]

平均摩擦系数：
\[
\bar C_f=\frac{1.328}{\sqrt{Re_L}}=0.0682.
\]

动压：
\[
\frac12\rho U^2=\frac12\times925\times0.6^2=166.5\ \mathrm{Pa}.
\]

平板两面面积：
\[
2S=2Lb=2\times0.5\times0.15=0.15\ \mathrm{m^2}.
\]

阻力：
\[
D=\bar C_f\frac12\rho U^2(2S)
=0.0682\times166.5\times0.15.
\]

\[
\boxed{\delta_{\max}=0.128\ \mathrm{m}}
\]

\[
\boxed{D=1.7\ \mathrm{N}}
\]

\section*{5. 书 10.3：排挤厚度和动量损失厚度}

\subsection*{题目}

假定层流边界层的速度分布如下，请计算边界层的排挤厚度 \(\delta_1\) 和动量损失厚度 \(\delta_2\)。

\[
(1)\quad u=U\frac{y}{\delta},\quad y<\delta;\qquad u=U,\quad y>\delta.
\]

\[
(2)\quad u=U\left[\frac{3}{2}\frac{y}{\delta}
-\frac{1}{2}\left(\frac{y}{\delta}\right)^3\right],
\quad y<\delta;\qquad u=U,\quad y>\delta.
\]

\subsection*{解答}

定义：
\[
\delta_1=\int_0^\infty\left(1-\frac{u}{U}\right)\dd y,\qquad
\delta_2=\int_0^\infty\frac{u}{U}\left(1-\frac{u}{U}\right)\dd y.
\]

令
\[
\eta=\frac{y}{\delta},\qquad \dd y=\delta\dd\eta.
\]

(1) \(\frac{u}{U}=\eta\)：
\[
\delta_1=\delta\int_0^1(1-\eta)\dd\eta=\frac{\delta}{2}.
\]

\[
\delta_2=\delta\int_0^1\eta(1-\eta)\dd\eta=\frac{\delta}{6}.
\]

(2) \(\frac{u}{U}=f=\frac32\eta-\frac12\eta^3\)：
\[
\delta_1=\delta\int_0^1(1-f)\dd\eta=\frac{3}{8}\delta.
\]

\[
\delta_2=\delta\int_0^1f(1-f)\dd\eta=\frac{39}{280}\delta.
\]

\[
\boxed{(1)\quad \delta_1=\frac{\delta}{2},\quad \delta_2=\frac{\delta}{6}}
\]

\[
\boxed{(2)\quad \delta_1=\frac38\delta,\quad \delta_2=\frac{39}{280}\delta}
\]

\section*{6. 书 10.5：正弦速度分布}

\subsection*{题目}

假设平板层流边界层内速度分布为
\[
\frac{u}{U_e}=a+b\sin\left(\frac{cy}{\delta}\right),
\]
式中 \(a,b,c\) 为待定常数，试用动量积分关系式方法求 \(\delta_1/\delta_2\) 和 \(C_D\)。

\subsection*{解答}

由边界条件：
\[
y=0:\quad u=0\Rightarrow a=0.
\]

\[
y=\delta:\quad u=U_e\Rightarrow b\sin c=1.
\]

\[
y=\delta:\quad \frac{\partial u}{\partial y}=0
\Rightarrow \cos c=0.
\]

故
\[
c=\frac{\pi}{2},\qquad b=1.
\]

速度分布为：
\[
\frac{u}{U_e}=\sin\left(\frac{\pi y}{2\delta}\right).
\]

令
\[
\eta=\frac{y}{\delta},\qquad f=\sin\frac{\pi\eta}{2}.
\]

\[
\delta_1=\delta\int_0^1(1-f)\dd\eta
=\delta\left(1-\frac{2}{\pi}\right).
\]

\[
\delta_2=\delta\int_0^1f(1-f)\dd\eta
=\delta\left(\frac{2}{\pi}-\frac12\right).
\]

\[
\boxed{\frac{\delta_1}{\delta_2}
=\frac{1-\frac{2}{\pi}}{\frac{2}{\pi}-\frac12}
\approx2.66}
\]

动量积分关系式：
\[
\frac{\tau_w}{\rho U_e^2}=\frac{\dd\delta_2}{\dd x}.
\]

壁面切应力：
\[
\tau_w=\mu\left(\frac{\partial u}{\partial y}\right)_0
=\mu U_e\frac{\pi}{2\delta}.
\]

设
\[
\delta_2=A\delta,\qquad A=\frac{2}{\pi}-\frac12.
\]

代入并积分得：
\[
C_D=\frac{2\delta_2(L)}{L}
=2\sqrt{\frac{A\pi}{Re_L}}.
\]

因为
\[
2\sqrt{A\pi}=1.31,
\]

\[
\boxed{C_D=\frac{1.31}{\sqrt{Re_L}}}
\]

\section*{7. 书 10.10：火车摩擦阻力功率}

\subsection*{题目}

流线型火车长 \(110\ \mathrm{m}\)、宽 \(2.75\ \mathrm{m}\)、高 \(2.75\ \mathrm{m}\)，假定顶部及所有侧面的表面摩擦力相当于长 \(110\ \mathrm{m}\)、宽 \(8.25\ \mathrm{m}\) 的光滑平板的一个表面。火车以 \(160\ \mathrm{km/h}\) 的速度在空气中等速行驶，空气密度
\[
\rho=1.22\ \mathrm{kg/m^3},
\]
粘度
\[
\mu=1.79\times10^{-5}\ \mathrm{Pa\cdot s},
\]
如临界 \(Re\) 为 \(5\times10^5\)，问层流边界层能维持多长？火车为克服风的摩擦阻力，消耗的功率为多少？

\subsection*{解答}

\[
U=\frac{160}{3.6}=44.44\ \mathrm{m/s}.
\]

等效平板面积：
\[
S=110\times8.25=907.5\ \mathrm{m^2}.
\]

层流边界层维持长度：
\[
x_{cr}=\frac{Re_{cr}\mu}{\rho U}
=\frac{5\times10^5\times1.79\times10^{-5}}{1.22\times44.44}.
\]

\[
\boxed{x_{cr}=0.165\ \mathrm{m}}
\]

全长雷诺数：
\[
Re_L=\frac{\rho UL}{\mu}
=\frac{1.22\times44.44\times110}{1.79\times10^{-5}}
=3.33\times10^8.
\]

转捩修正平均摩擦系数：
\[
\bar C_f=\frac{0.074}{Re_L^{1/5}}-\frac{1742}{Re_L}
=0.001456.
\]

摩擦阻力：
\[
D=\bar C_f\frac12\rho U^2S
=0.001456\times\frac12\times1.22\times44.44^2\times907.5.
\]

\[
D=1.59\times10^3\ \mathrm{N}.
\]

功率：
\[
P=DU=1.59\times10^3\times44.44.
\]

\[
\boxed{P\approx70.8\ \mathrm{kW}}
\]

\section*{8. 平板单面摩擦阻力和尾部边界层厚度}

\subsection*{题目}

已知来流 \(2.5\ \mathrm{m/s}\) 绕过零攻角平板，平板长 \(L=50\ \mathrm{cm}\)，宽 \(b=3\ \mathrm{m}\)，计算并比较不同流体绕流下，平板单面摩擦阻力及平板尾部的边界层名义厚度。临界
\[
Re_{cr}=5\times10^5.
\]

(a) 空气：
\[
\rho=1.2\ \mathrm{kg/m^3},\qquad \nu=1.5\times10^{-5}\ \mathrm{m^2/s}.
\]

(b) 水：
\[
\rho=998\ \mathrm{kg/m^3},\qquad \nu=1.005\times10^{-6}\ \mathrm{m^2/s}.
\]

\subsection*{解答}

\[
U=2.5\ \mathrm{m/s},\qquad L=0.5\ \mathrm{m},\qquad b=3\ \mathrm{m}.
\]

\[
S=Lb=1.5\ \mathrm{m^2}.
\]

\subsubsection*{(a) 空气}

\[
Re_L=\frac{UL}{\nu}
=\frac{2.5\times0.5}{1.5\times10^{-5}}
=8.33\times10^4.
\]

\[
Re_L<5\times10^5,
\]
所以全层流。

\[
\bar C_f=\frac{1.328}{\sqrt{Re_L}}
=0.00460.
\]

\[
D=\bar C_f\frac12\rho U^2S
=0.00460\times\frac12\times1.2\times2.5^2\times1.5.
\]

\[
\boxed{D_{\mathrm{air}}\approx0.0259\ \mathrm{N}}
\]

\[
\delta_L=\frac{5L}{\sqrt{Re_L}}
=\frac{5\times0.5}{\sqrt{8.33\times10^4}}
=8.66\times10^{-3}\ \mathrm{m}.
\]

\[
\boxed{\delta_{\mathrm{air}}\approx8.66\ \mathrm{mm}}
\]

\subsubsection*{(b) 水}

\[
Re_L=\frac{UL}{\nu}
=\frac{2.5\times0.5}{1.005\times10^{-6}}
=1.24\times10^6.
\]

\[
Re_L>5\times10^5,
\]
发生转捩。

转捩位置：
\[
x_{cr}=\frac{Re_{cr}\nu}{U}
=\frac{5\times10^5\times1.005\times10^{-6}}{2.5}
=0.201\ \mathrm{m}.
\]

平均摩擦系数：
\[
\bar C_f=\frac{0.074}{Re_L^{1/5}}-\frac{1742}{Re_L}
=0.00307.
\]

\[
D=\bar C_f\frac12\rho U^2S
=0.00307\times\frac12\times998\times2.5^2\times1.5.
\]

\[
\boxed{D_{\mathrm{water}}\approx14.4\ \mathrm{N}}
\]

尾部为湍流边界层，名义厚度：
\[
\delta_L=\frac{0.37L}{Re_L^{1/5}}
=\frac{0.37\times0.5}{(1.24\times10^6)^{1/5}}.
\]

\[
\boxed{\delta_{\mathrm{water}}\approx11.2\ \mathrm{mm}}
\]

\begin{center}
\begin{tabular}{lllll}
\toprule
流体 & \(Re_L\) & 状态 & 单面摩擦阻力 & 尾部边界层厚度\\
\midrule
空气 & \(8.33\times10^4\) & 全层流 & \(0.0259\ \mathrm{N}\) & \(8.66\ \mathrm{mm}\)\\
水 & \(1.24\times10^6\) & 转捩后湍流 & \(14.4\ \mathrm{N}\) & \(11.2\ \mathrm{mm}\)\\
\bottomrule
\end{tabular}
\end{center}

\end{document}
